\section{Resultados y discusión}

\subsection{Nanopartículas obtenidas}

% Expandir un poco más en que cambia entre las 3 y los colores
Se realizaron tres síntesis diferentes de AgNPs. La primera fue de los  nanotriángulos (AgTNPs) y fue la única que se pudo caracterizar parcialmente ya que los microscopios sufrieron fallas de funcionamiento y se encuentran en reparación. Las otras dos síntesis de partículas esféricas de distinto tamaño no llegaron a ser caracterizadas.

La figura \ref{fig:campo_claro} muestra imágenes de campo claro representativas de las AgTNPs sintetizadas, utilizando el microscopio Phillips CM200-UT. La primer característica que se destaca es la superposición de las partículas y su orientación preferencial. Se identifican claramente que dicha superposición por que son parcialmente transparente a los electrones. Esto es un indicador indirecto del delgado espesor de estos AgTNps. Se puede observar un amplio rango de tamaños, con lados entre \qtyrange{20}{120}{\nano\meter}. 

% TODO: Acomodar figuras y escribir epígrafe
% Me combiene juntarlas en algún editor de imágen y ponerles (a) (b) a mano para controlar más la distribución?
\begin{figure}[h]
	\centering
	\begin{subfigure}{\linewidth}
		\centering
		\includegraphics[width=0.9\linewidth]{./NPs/mias/012.png}
		\caption{}
		\label{fig:cumulo_NPs}
	\end{subfigure}
	
	\begin{subfigure}{\linewidth}
		\centering
		\includegraphics[width=0.9\linewidth]{./NPs/mias/03.png}
		\caption{}
		\label{fig:pocas_NPs}
	\end{subfigure}
	
	\label{fig:campo_claro}
	\caption{Imágenes de campo claro de AgNPs triangulares azules tomadas utilizando el CM200. Se observa conglomeración moderada y gran variedad en el largo de los lados.}
\end{figure}

Nótese que en la figura \ref{fig:pocas_NPs} se observan unas franjas oscuras y claras en la superposición de algunas AgTNPs. Esos patrones de Moirè son producto de la superposición de dos estructuras periódicas levemente rotadas o desfasadas. No se deben confundir con espaciados interplanaras que se observan el HRTEM.


\subsubsection*{Mecanismo de formación de los AgTNPs}\label{SEC:mecanismo_formacion_particulas_triangulares}

\begin{scheme}
	\centering \scriptsize \sffamily
	\includesvg[width=\linewidth]{NPs/reaccion}
	\caption{Formación de AgTNPs. En la primer etapa se forman partículas muy pequeñas que actúan como semillas y luego crecen adoptando formas triangulares o hexagonales. El \ce{NaBH4} reduce los iones plata \ce{A+} y comienza a formar semillas. El \ce{H2O2} oxida las partículas de plata.} 
	\label{sch:reaccion}
\end{scheme}

% Rol de cada reactivo
En términos generales, la síntesis de este tipo de NPs corresponde a la formación de un coloide estabilizado electrostáticamente. En el esquema \ref{sch:reaccion} se describe el mecanismo aceptado para la formación de estas AgTNPs. El \ce{AgNO3} es el precursor metálico, que aporta los iones \ce{Ag+} en solución.
Los aniones citrato (\ce{C6H5O7^3-}) actúan como agente protector o estabilizador. 
Se adsorbe a la superficie de las partículas durante el crecimiento y aporta carga negativa.
De esta manera se evita que las AgNPs en formación se aglomeren, coalescan y precipiten como un sólido masivo.
El \ce{BH4-} es un agente reductor fuerte, que reduce los iones \ce{Ag+} a \ce{Ag^0}. 
El reductor se encuentra en exceso (10:1 respecto de \ce{Ag+}) dentro de la solución, por lo que se forma una gran cantidad de núcleos de \ce{Ag^0} que no llegan a crecer a tamaños mayores a los $~$\qty{3}{\nano\meter}.
El peróxido de hidrógeno (\ce{H2O2}) aquí actúa como oxidante de las AgTNPs recién formadas (\textit{etching}). Así se inicia una reacción cíclica de reducción/oxidación que lleva a la formación de nanoplatos. 

\begin{figure}
	\centering \sffamily \tiny
	\includesvg[width=0.9\linewidth]{NPs/forma}
	\caption{Diagrama de las caras cristalinas presentes en los nanoplatos de Ag.}
	\label{FIG:forma}
\end{figure}

Se ha propuesto que este mecanismo cíclico reducción producida por el \ce{NaBH4} y la oxidación producida por el \ce{H2O2} fomentan la formación de fallas de apilamiento el durante el crecimiento de las nanopartículas. Rocha y colaboradores mostraron que las AgNPs semillas esféricas presentan estas fallas de apilamiento en los planos \hkl{111}.\cite{rocha2007} Así entonces se pueden describir los nanoplatos como estructuras cristalinas fcc, donde las fallas de apilamiento generan una capa de estructura hcp, figura \ref{FIG:forma}. En la figura \ref{fig:hrtem_canto} se muestra una imagen experimental de una AgTNPs en una vista lateral (de canto) en un eje de zona cercano a \hkl[110] donde se evidencian las distintas fallas de apilamiento y las regiones con estructura fcc-hcp-fcc. Recientemente Tan y colaboradores observaron una disposición de distintos tipos de fallas en los planos basales \hkl{111}.\cite{tan2021}

% TODO: Explicar mejor la justificación del rol de las fallas
%La presencia de fallas causa que a las partículas se agreguen preferencialmente sobre estas zonas.
%Entonces, en lugar de crecer de manera isotrópica, lo hace de manera anisotrópica formando platos. 
%De la misma manera, ciertos planos laterales se ven favorecidos, de esta manera, en lugar de crecer platos hexagonales de se forman platos triangulares de plata.
Hasta el momento, el modelo que permite relacionar la geometría que adoptan las partículas se basa en la presencia de las fallas de apilamiento. Asimismo, el hecho que los iones citrato se unan fuertemente a las caras \hkl{111} acompaña la justificación de la geometría adoptada durante el crecimiento. Se asume que los iones citrato al adsorberse sobre las caras \hkl{111} las "bloquea" (pasivasión) y los nuevos átomos de \ce{Ag^0} se adsorben preferentemente sobre las caras \hkl{100} o las caras con múltiples fallas. Sin embargo, este modelo es insuficiente para explicar el alto grado de homogeneidad en el estrecho espesor de estas partículas. 


\subsection{HRTEM - Experimental y Simulaciones}



% Aparecen dos puntos es demasiado coloquial


% Intento de HRTEM
\begin{figure}
	\centering
	\includegraphics[width=\linewidth]{./NPs/mias/014_mas_FFT.png}
	\label{fig:mi_HRTEM}	
	\caption{Imágen de alta resolución del borde de dos partículas superpuestas. En la esquina superior izquierda se muestra la transformada de Fourier de la imagen (FFT). Los dos puntos presentes en la FFT indican la presencia de dos diferentes orientaciones de los planos cristalinos. La forma oblonga está asociada a astigmatismo no corregido en las lentes. }
\end{figure}
% Por qué no se observan el resto de los puntos en la FFT?

\begin{figure}
	\centering
	\includegraphics[width=\linewidth]{./NPs/mias/015_con_FFT.png}
	\label{fig:mi_HRTEM_2}	
	\caption{Imágen de alta resolución de una partículas. En la esquina superior derecha se muestra la transformada de Fourier de la imagen (FFT). Las reflexiones corresponden a una estructura hcp.
	}
\end{figure}


Luego se realizaron imágenes de HRTEM de la misma muestra. En Fig. \ref{fig:mi_HRTEM} se tomaron dos partículas apiladas. A pesar del alto nivel de ruido y que no se observen columnas atómicas, igualmente aparecen dos puntos en la transformada de Fourier. Estos están asociados a los planos \hkl{} que se encuentran paralelos a la dirección del borde. Los otros planos no se observan por el astigmatismo de los lentes.

En Fig. \ref{fig:mi_HRTEM_2} se observa una partícula más pequeña cuyo FFT presenta 6 puntos en forma de hexágono. Esto indica la presencia de una estructura hcp en la partícula, lo que son las fallas de apilamiento. A diferencia de la figura \ref{fig:mi_HRTEM} la FFT no presenta tanta deformación en la diagonal. Es decir, no presenta astigmatismo. Esta diferencia permite que aparezcan las 6 reflexiones esperadas y no solo 2 que coinciden con la dirección en que se enlonga la FFT.

La periodicidad los planos identificados en ambas figuras son de \SI{0.26 \pm 0.1}{\nano\metre}, en concordancia con los \SI{0.249}{\nano\metre} reportados en la literatura \cite{rocha2007}.
% Están realmente en cooncordancia la imagen del borde da una distancia más cercana a 0.27
% Y la segunda da 0.26, me tiraria a decir que mantenemos el problema de calibración pero no tengo
% suficiente seguridad para plantearlo.

% Agregar la imagen de la otra partícul que conseguimos
% Tiene menos atigmatismo entonces se ven las 6 reflexiones aunque no sean obvias en la imagen


%TODO: Reemplazar por figura original o hablar de la figura que tomó Julio
%TODO: Marcar las posiciones de las fallas
\begin{figure}
	\centering
	\includegraphics[width=\linewidth]{NPs/Julio/03-HRTEMcanto.png}
	\caption{Imagen de una nanopartícula triangular de plata vista desde el costado. El haz de electrones da con el plano {hkl} }
	\label{fig:hrtem_canto}
\end{figure}

%TODO Discutir la imgen anterior

En Fig. \ref{fig:hrtem_canto} se observan las posiciones de distintas fallas a lo largo de la nanopartícula.

%TODO: Comparar con Tan2021

Luego se analizó la imagen presente en \ref{fig:NP_grande}. Se observan varios fenónemos. 
Lo primero es que el contraste no es el mismo a lo largo de toda la partícula. 
En los bordes hay puntos negros sobre fondo blanco y en el centro se invierte.
E

\begin{figure}
	\centering
	\includegraphics[width=\linewidth]{./NPs/Julio/Julio_0022_full.png}
	\caption{Micrografía tomada mediante HRTEM de una nanopartícula de plata triangular. Hay una diferencia de contraste entre los bordes de la partícula y el  centro.}
	\label{fig:NP_grande}
\end{figure}

\begin{figure}
	\centering
	\includegraphics[width=\linewidth]{./NPs/Julio/Julio_0022_zoom_sin_escala.png}
	\caption{Aumento de la zona superior izquierda de Fig. \ref{fig:NP_grande}.}
	\label{fig:NP_grande_zoom}
\end{figure}

Lo primero se obtuvo una imagen simulada con el mismo tipo de contraste que en el borde nanopartícula.
A partir de estas se compararon los perfiles de intensidad en las tres direcciones de crecimiento (direcciones de crecimiento? ) de la nanopartícula.
A partir de los perfiles obtenemos cual es la periodicidad del patrón observado en la imagen.


%TODO: En la parte de la toma de perfiles, no me van a preguntar directamente por qué no trabajo en el espacio recíproco?
%TODO: Análsis de las simulaciones
% con la posición de los puntos...
%Para la imagen experimental se observa una distancia de \SI{Valor}{unidad}
%En el caso de la imagen simulada es \SI{Valor}{unidad}

%TODO: Figura de una simulación
%TODO: Comparación con el resultado experimental